\section{Introducción}

% The figure of 150.6 km² represents the raw total land area of the entire municipality,
% not just the urbanized footprint. Because of this, the average figure of roughly 1,385 inhabitants/km² is indeed a massive underestimation of the actual population density
% in the areas where people actually live.
% Here is how the territorial reality and the legal framework explain this discrepancy:
% 1. The Reality of the Terrain
% A vast portion of the municipality of Santa Cruz de Tenerife is composed of the Macizo de Anaga. This area, known as the Parque Rural de Anaga, is a declared Biosphere Reserve and features steep, mountainous terrain that is largely uninhabitable. Because this massive natural expanse is legally part of the municipality's total surface area, dividing the total population by the total km² artificially dilutes the density metric. If you were to look exclusively at the urbanized footprint (the city center, the coastal strip, and the southwest neighborhoods), the population density would be drastically higher, resembling that of major dense European cities.   2. Legal Protection and Land ClassificationThe reason this steep mountainous area remains uninhabited is not just topographical, but strictly regulated by law. If you reference the file named "Ley del Suelo, Planes de Ordenación Urbana", you will see that territorial and urban planning instruments mandate the strict delimitation and protection of these spaces.According to the law, island and municipal planning must identify and classify specific zones of the territory that are to be excluded from development, including:Areas that must be preserved from the urbanization and building process because transforming them would be incompatible with sustainable development.   Areas that must be preserved due to their existing or potential landscape, agricultural, or heritage value.   Areas that must be excluded from the urbanization process to prevent risks—such as geological, seismic, or meteorological risks (which is highly relevant for steep ravines and mountainsides).   Therefore, under the "Ley del Suelo, Planes de Ordenación Urbana", the mountain ranges of Anaga are strictly classified as protected rural land (suelo rústico / espacios naturales protegidos), effectively fencing them off from urban expansion.

El número de 150,6~km\textsuperscript{2} representa la superficie territorial total del término municipal de Santa Cruz de Tenerife, no únicamente la superficie urbanizada. Por esta razón, el indicador promedio de aproximadamente 1.385~hab/km\textsuperscript{2} supone una subestimación sustancial de la densidad demográfica real soportada en los sectores donde efectivamente reside la población.

Esta divergencia entre la cifra agregada y la realidad territorial se explica a través de dos condicionantes fundamentales:

\begin{enumerate}
  \item \textbf{La realidad física y orográfica del relieve:} Una vasta extensión del término municipal está integrada por el Macizo de Anaga. Este ámbito, protegido bajo la figura de Parque Rural y declarado Reserva de la Biosfera por la UNESCO, presenta una topografía abrupta, montañosa y de acusadas pendientes que resulta materialmente inhabitable y hostil al desarrollo urbano. Dado que este inmenso espacio natural forma parte administrativamente de la superficie total del municipio, el cálculo aritmético que divide la población total entre el área municipal diluye de manera artificial el índice de densidad. Si el análisis se circunscribe a la huella urbanizada real (el centro histórico y ensanches, la fachada litoral y las barriadas del suroeste), la densidad de población neta se incrementa de forma notable, alcanzando registros equiparables a los de las grandes y densas metrópolis europeas.
  \item \textbf{Protección jurídica y clasificación del suelo:} La ausencia de desarrollo urbano en las zonas montañosas no responde únicamente a barreras orográficas, sino a una estricta tutela jurídica establecida en la legislación del suelo y los instrumentos de ordenación territorial y urbanística. La normativa exige que los planes de ordenación identifiquen, delimiten y preserven preceptivamente aquellos ámbitos del territorio excluidos de todo proceso de desarrollo, entre los que destacan:
  \begin{itemize}
    \item Los terrenos que deban preservarse de la transformación urbanística y edificatoria por resultar esta incompatible con el desarrollo sostenible.
    \item Las áreas que deban salvaguardarse en función de sus valores paisajísticos, agrícolas, forestales, ecológicos, naturales o patrimoniales.
    \item Las zonas que deban excluirse de la urbanización con el fin expreso de prevenir y mitigar riesgos naturales y geotécnicos severos (riesgos geológicos, sísmicos o meteorológicos, especialmente críticos en laderas escarpadas y cauces de barrancos).
  \end{itemize}
  En aplicación de este régimen urbanístico, las estribaciones y crestas del Macizo de Anaga están clasificadas rigurosamente como Suelo Rústico Protegido (integrado en la Red Canaria de Espacios Naturales Protegidos), conformando una barrera geográfica y jurídica infranqueable frente a la expansión urbana del municipio.
\end{enumerate}

Bajo estos condicionantes de constricción física y delimitación legal del espacio habitable, se analizan a continuación las distintas fuerzas demográficas, socioeconómicas y normativas que actúan sobre el urbanismo y la ordenación del territorio en Santa Cruz de Tenerife desde 2011 hasta la actualidad.

\subsection{Turismo y Percepción: La Disparidad entre Alojamiento y Afluencia}

% 1. The 2011–2012 Population Drop: Factual Exodus or Statistical Artifact?
% It was purely an administrative and statistical artifact, not a physical exodus of 15,000 residents.

% What actually happened:
% The 2011 Census Audit: In 2011–2012, the INE conducted Spain's decennial Censo de Población y Viviendas 2011, transitioning from door-to-door paper forms to cross-referencing administrative databases (Padrón, Seguridad Social, and the Ministry of the Interior's National DNI/NIE database).
% The "Habitantes Fantasma" of Santa Cruz: When the INE audited Santa Cruz de Tenerife's municipal register, it discovered an accumulated backlog of invalid, duplicated, or non-existent DNI entries that had been carried over since 1996.
% The Official Purge: The INE rejected the city's official submission and forced an administrative cleanup of 15,306 entries on 1 January 2012 (a sudden drop from 222,271 down to 206,965, or -6.89%).
% The Political & Fiscal Scandal: Because municipal funding allocations (such as the Fondo Canario de Financiación Municipal, state tax shares, and Canary REF revenue) were distributed per capita based on padrón figures, this purge caused a major municipal and political controversy across the Canary Islands.

% 2. Would plotting these series showcase the need for efficient land use management?
% Yes, absolutely. In fact, it provides a much more sophisticated urban planning argument than the typical "runaway demographic growth" narrative.

% Here is why and how it builds your thesis:

% A. The Paradox of "Flat Growth" vs. "Severe Urban Tension"
% Looking at the data extracted from your files:

% Resident Population (population_all.json): Post-2012, population has been virtually flat (from 206,965 in 2012 to 211,498 in 2025, barely +2.18% over 13 years).
% Touristic Lodging Population (alojatur_poblacion_tur_etc.json): Santa Cruz is the provincial administrative and port capital, not a resort enclave like Arona or Adeje. Its regulated tourist population is modest (oscillating between ~900 and ~1,700 daily equivalent tourists), though pressure from vivienda vacacional (VV) and daytime excursionists in central quarters has surged recently.
% Household Income (renta.json): Average gross income rose steadily from 26,124 € (2013) to 28,662 € (2021).
% B. The Planning Thesis (The Supply Bottleneck)
% In conventional planning debates, housing shortages and spatial stress are blamed on explosive demographic booms. In Santa Cruz, the data shows the opposite:

% Shrinking Household Size: Even with flat population, households have unbundled (more single-person and two-person homes), meaning the city requires significantly more housing units just to house the same number of people.
% Economic & Commercial Vitality: Rising income drives demand for better housing quality, logistics space, and tertiary services.
% The Artificial Bottleneck (Planning Paralysis): Because the PGO 2013 was annulled by the TSJC (Sentence 116/2014) and the 2022 provisional ordinance was suspended (Sentence 70/2023), the city is legally anchored to the obsolete 2005 plan.
% Land transformation in Sectorized Urbanizable Land (SUSO) and Unconsolidated Urban Land (SUNC) is virtually frozen.
% Private and public developers cannot deliver new housing stock or complete strategic urban stitching (e.g., around the Barranco de Santos or Cabo Llanos).

Desde el comienzo de los datos disponibles (2011-2025), la poblacion de Santa Cruz de Tenerife se ha mantenido estable entre 205,000 y 210,000 habitantes.
% , pero la composicion de los hogares ha cambiado, con mas hogares unipersonales y hogares de dos personas, lo que requiere mas viviendas para albergar al mismo numero de personas. Ademas, la renta ha aumentado, lo que impulsa la demanda de viviendas de mayor calidad, espacios logisticos y servicios terciarios.
% Sin embargo, el Plan General de Ordenacion de 2013 fue anulado por el Tribunal Superior de Justicia de Canarias (sentencia 116/2014) y la ordenanza provisional de 2022 fue suspendida por el mismo tribunal (sentencia 70/2023), lo que ha congelado la transformacion de suelos urbanizables y no consolidados, impidiendo que los promotores privados y publicos entreguen nuevas viviendas o completen proyectos urbanos estrategicos.
La poblacion turistica recogida por el ISTAC muestra que en términos de alojamiento regulado apenas alcanza los 2,000 habitantes.
Aunque ha aumentado en un 50\% en el periodo considerado, una alta crecimiento de una cifra pequena todavia produce una cifra pequena.

\begin{figure}[htbp]
  \centering
  \includegraphics[width=0.95\linewidth]{figures/demographics_trend.pdf}
  \caption{Evolución de la Población Residente y Turística en Santa Cruz de Tenerife (2011-2025). Fuente: Elaboración propia a partir de datos del ISTAC e INE.}
  \label{fig:demografia_trend}
\end{figure}

%  there is an enormous disparity between the perceived tourist presence in Santa Cruz de Tenerife and the actual resident lodging demand, driven almost entirely by day-visitors (cruceristas and excursionists).

% Here is the breakdown based on official data from the Autoridad Portuaria de Santa Cruz de Tenerife (Puertos de Tenerife), Aena, and the ISTAC:

% 1. The Numbers: Cruise Passengers vs. Air Tourists
% Category  Volume (Annual)  Modality / Behavior in Santa Cruz  Impact on Housing Stock
% Cruise Tourists (Puerto de S/C)  ~972.000 (2024)
% (817.647 en 2023)  Day-trippers (Transeúntes de escala):
% • ~65% (~630.000) disembark for 4–8 hours in the city center and sleep on the ship.
% • ~35% (~340.000) are puerto base (start/end itinerary).  Zero housing impact (they sleep on the ship; they only consume sidewalks, terraces, and retail).
% Air Tourists (Tenerife Island Total)  ~6.500.000 (2023–2024)
% (TFS + TFN)  Concentrated in resort enclaves:
% • 75%–80% stay in the South (Adeje, Arona, Guía de Isora).
% • 15% in the North (Puerto de la Cruz).
% • Barely 2%–3% stay overnight in Santa Cruz.  Localized pressure in southern municipalities, not Santa Cruz.
% Regulated Overnight Tourists in Santa Cruz  ~1.500 – 1.700
% (Población equivalente diaria, Alojatur)  Urban / business / cultural hotel lodging with an average stay of 2–3 nights.  Marginal (<0.8% of the municipal resident population).
% Day Excursionists from the South/North  Hundreds of thousands  Tourists staying in Playa de las Américas, Adeje, or Puerto de la Cruz visiting Santa Cruz for shopping, museums, or Anaga.  Zero housing impact (they return to their southern hotels before dinner).
% 2. Why Perception Disconnects from Reality
% Spatial Concentration in the Historic Core: Cruise ships dock directly at the Muelle Sur and Muelle de Ribera, just steps from Plaza de España, Calle Castillo, Plaza de la Candelaria, and the Mercado de Nuestra Señora de África.
% On peak winter days (October–April), when 2 or 3 large liners dock simultaneously, 8,000 to 12,000 passengers flood a narrow, pedestrianized central corridor between 09:00 and 17:00.
% This creates a vivid visual impression of intense "over-tourism", even though not a single one of these visitors rents an apartment or stays in a local hotel bed.
% Daytime Excursionists: Tourists vacationing in the South take rental cars or TITSA express buses (Línea 110/111) to spend the day in the capital, shopping or touring the Anaga Biosphere Reserve, concentrating human traffic in central streets without registering in municipal lodging statistics.
% 3. Urban Planning & Housing Takeaway (For your Introduction)
% This distinction sharpens your argument in

% sections/15_introduccion.tex
% :

% The Fallacy: Attributing the housing shortage and spatial tension in Santa Cruz de Tenerife to "uncontrolled tourist lodging."
% The Reality:
% Nearly 1 million cruise passengers per year create high pedestrian pressure on public space, street cleaning, and retail, but do not compete for residential housing.
% Over-night lodging in the capital remains tiny (~1,600 equivalent residents).
% Therefore, the housing affordability crisis is almost entirely an internal supply crisis:
% Demographic unbundling: Flat population, but smaller households requiring more individual units.
% Regulatory paralysis: An annulled PGO (2013) and suspended provisional ordinance (2022) that have frozen the legal transformation of new residential sectors (SUSO) and urban regeneration units (SUNC).

A esta baja presión de alojamiento reglado se contrapone, sin embargo, una fuerte percepción ciudadana de saturación turística.
Dicho fenómeno responde fundamentalmente al turismo de día:
el Puerto de Santa Cruz de Tenerife supera los 900.000 cruceristas anuales (con más de 600.000 transeúntes de escala que desembarcan por unas horas en el casco histórico),
a los que se suman las excursiones diurnas procedentes de los núcleos turísticos del sur.
Este flujo flotante intensifica el uso del espacio público y los servicios urbanos durante el día, pero no consume parque habitacional residencial,
confirmando que la escasez de vivienda en el municipio no obedece a una absorción masiva por turismo alojativo.
%, sino a la parálisis estructural de la gestión del suelo.

\subsection{Inmigración}

% 1. The Latin American Inversion (Historical Flow to America Reversed)
% Status: 100% Factual and Accurate.
% The Historical Context: For over a century, the Canary Islands were an emigration society. Tens of thousands of Canarians emigrated to America—chiefly Cuba (early 20th century) and Venezuela (1950s–1980s, historically known as "la octava isla").
% The Modern Reversal:
% In the last 15 years, that flow has completely inverted. Due to political and economic instability in Latin America (notably Venezuela, Colombia, and Cuba), hundreds of thousands of people have migrated to the Canary Islands.
% A substantial portion are returnees and descendants (children and grandchildren of Canarians) who already hold or easily acquire Spanish nationality under the Ley de Memoria Democrática (Ley de Nietos).
% In Santa Cruz de Tenerife, registered foreign nationals alone rose from 4,473 in 2000 to a record 16,776 in 2025 (ISTAC), not including the tens of thousands who hold Spanish passports by descent.
% Impact on Housing Demand: Very High and Direct.
% These individuals and families settle permanently as long-term residents (vivienda habitual).
% They enter the formal local labor market (healthcare, services, commerce, education) and compete directly in the residential rental and purchase market across Santa Cruz (Ofra, Añaza, La Salud, Tome Cano, Centro).
% 2. Irregular Migration from Africa (Cayucos / Pateras)
% Status: Crucial Correction — It does NOT significantly impact residential housing demand.
% Why it does not affect the private housing market:
% A Transit Phenomenon, Not Settlement: The Ruta Canaria is an emergency transit border route. Over 90% to 95% of adult irregular migrants arriving by sea do not settle in the Canary Islands. The Ministry of Inclusion systematically organizes humanitarian transfers (derivaciones) to mainland Spain (la Península) and onwards to Europe.
% Institutional Reception Centers (No Private Rental Access): During their temporary stay in the islands, they are housed in state-run emergency facilities (such as Las Raíces in La Laguna or Canarias 50 in Las Palmas) managed by NGOs (Cruz Roja, CEAR, Accem).
% Lack of Documentation: Without Spanish DNI/NIE, valid work permits, and employment contracts (nóminas), it is virtually impossible to sign private residential rental leases or secure mortgages.
% Unaccompanied Minors (Menores Tutelados): Those who remain under the custody of the Gobierno de Canarias are housed in public youth shelters and emergency care homes, not in the private real estate market.
% Urban Takeaway: Conflating the humanitarian crisis in emergency reception centers with the residential housing shortage is a common misconception. The housing shortage is driven by permanent residents, household unbundling, and planning paralysis.
% 3. The Pope's Visit
% Status: Factual in its humanitarian focus, but unrelated to local housing policy.
% The Context:
% Pope Francis explicitly announced his desire to visit the Canary Islands (reiterated in meetings with Canarian institutions) to spotlight the humanitarian drama of the Atlantic route, comparing it to his iconic 2013 visit to Lampedusa.
% His message focused strictly on human dignity, preventing the Atlantic from becoming a "cemetery without tombstones," and commending the humanitarian solidarity of the Canarian people.
% It was a moral and pastoral mission focused on sea rescues and human rights, not an urban-planning or housing-market critique.
% Summary: Demographic Drivers vs. Housing Impact in Santa Cruz de Tenerife
% Demographic / Visitor Flow  Volume  Nature of Stay  Impact on Santa Cruz Housing Stock
% Latin American & EU Migration  High & growing (Record foreign registration + Spanish descendants)  Permanent settlement (family homes, long-term employment)  Direct & High: Drives primary residential rental & purchase demand.
% Cruise Passengers  ~972.000 / year  Day-trippers (4–8 hours)  Zero: Sleep on ships; only impact street traffic & retail.
% Air Tourists (Canary total)  ~6.5M / year  Resorts (South/North) (Only ~1,600 stay in S/C hotels daily)  Marginal in S/C (except localized vivienda vacacional pressure).
% Irregular Sea Migration (Cayucos)  Humanitarian emergency  Transit (State reception facilities ➔ mainland transfers)  Negligible: Does not enter the private residential rental market.
% Internal Demographic Shift  Flat total population  Household Unbundling (more 1-2 person homes)  Very High: Same population requires thousands more dwellings.

% ==============================================================================
% ANÁLISIS DETALLADO: DINÁMICA DE FLUJOS MIGRATORIOS Y SALDO EXTERIOR (INE / ISTAC)
% ==============================================================================
% 1. Panel Superior: Entradas vs. Salidas Exteriores (Flujos Transatlánticos)
% - La crisis financiera española (2008–2014):
%   Las entradas desde el extranjero cayeron a su mínimo (~10.300 personas/año en 2009–2010).
%   Las salidas de residentes hacia el extranjero se duplicaron hasta alcanzar un máximo de 8.867 salidas en 2015 (salida de jóvenes y retorno temporal de extranjeros durante la recesión).
% - La inversión y ola migratoria de América Latina (2015–actualidad):
%   La profundización de la crisis institucional y económica en Venezuela, Colombia y Cuba provocó un giro histórico: las llegadas procedentes del exterior crecieron vertiginosamente de 14.199 (2015) a más de 32.800 personas/año (2022–2024).
%   Por el contrario, la emigración exterior se redujo notablemente, abriendo una brecha neta masiva a favor de las entradas.
%
% 2. Panel Inferior: Saldo Migratorio Neto Municipal (Santa Cruz de Tenerife)
% A escala estrictamente municipal (término de Santa Cruz de Tenerife, código 38038), los datos del ISTAC e INE muestran dos regímenes diametralmente opuestos:
% - Fase expulsora / pérdida neta (2009–2015):
%   Saldos fuertemente negativos: -2.168 hab. en 2010, -945 hab. en 2011 y -1.618 hab. en 2014. Santa Cruz perdía población por el freno inmigratorio y la salida hacia el exterior u otros municipios.
% - Fase de inversión y presión residencial sostenida (2016–2024):
%   A partir de 2016, el saldo municipal se vuelve intensamente positivo: +2.014 (2017), +2.545 (2018), +2.457 (2019), +2.084 (2022), +2.618 (2023) y +2.082 (2024).
%   Este superávit neto anual continuo de más de 2.000 nuevos residentes netos cada año en la capital explica directamente por qué la demanda habitacional permanente presiona el mercado de compra y alquiler, a pesar de que el padrón global total fluctúe de forma más moderada por el envejecimiento vegetativo de la población autóctona.
%
% 3. El trasfondo histórico: De "la Octava Isla" al retorno generacional
% - Flujo histórico (siglos XIX a 1980): Canarias fue durante más de un siglo una sociedad de emigración masiva hacia América (primero Cuba en los años 20 y luego Venezuela en los 50–70, popularizada como "la octava isla").
% - Flujo contemporáneo invertido: La inestabilidad en el continente americano trajo de vuelta a emigrantes retornados, a sus hijos y nietos (acogiéndose a la Ley de Memoria Democrática), así como a nuevos contingentes latinoamericanos.
% - Esto explica que el registro de extranjeros censados en Santa Cruz pasara de 4.473 en el año 2000 a 16.776 en 2025 (casi el cuádruple), a los que se suman decenas de miles que acceden directamente con pasaporte español por descendencia.
% ==============================================================================

En el plano demográfico, la presión sobre el parque de viviendas tampoco obedece a los flujos humanitarios de la ruta subsahariana, un fenómeno de tránsito gestionado en centros de acogida e institucionales sin incidencia en el mercado de alquiler privado.

Por el contrario, la demanda habitacional proviene de la inversión del histórico flujo migratorio con América:
Santa Cruz ha pasado de ser un foco de emigración a un receptor de población residente (especialmente retornados y ciudadanos de origen latinoamericano con arraigo o nacionalidad).
%, que sumado al encogimiento del tamaño medio de los hogares, incrementa la necesidad objetiva de viviendas permanentes en un mercado bloqueado por la falta de suelo de desarrollo.

El registro de extranjeros en Santa Cruz de Tenerife pasó de 4.473 personas en el año 2000 a un récord de 16.776 en 2025, sin incluir a los miles que ya cuentan con pasaporte español por vínculos familiares.

\begin{figure}[htbp]
  \centering
  \includegraphics[width=0.95\linewidth]{figures/migration_trend.pdf}
  \caption{Evolución de los Flujos Migratorios Exteriores y Saldo Migratorio Neto (2008--2024). Fuente: Elaboración propia a partir de datos de la Estadística de Migraciones y Cambios de Residencia (EMCR/INE) y la Estadística de Variaciones Residenciales (EVR/ISTAC).}
  \label{fig:migracion_trend}
\end{figure}

Como reflejan las series oficiales del INE y del ISTAC (Figura~\ref{fig:migracion_trend}), tras el ciclo de recesión 2008--2014 —caracterizado por la contracción de entradas y saldos netos negativos de hasta $-2.168$ habitantes en la capital—, el flujo exterior experimentó una inversión radical a partir de 2016. Las entradas procedentes del extranjero superan actualmente las 30.000 personas anuales a escala provincial, consolidando en el municipio de Santa Cruz de Tenerife un saldo migratorio neto positivo ininterrumpido superior a los 2.000 residentes permanentes al año ($+2.618$ en 2023). Esta afluencia continuada de población con vocación de asentamiento definitivo genera una presión directa sobre el parque de vivienda residencial, en un mercado constreñido por la falta de suelo de desarrollo.

\subsection{Familias, Nuevos Hogares y Exigencias Habitacionales}

% The Three Structural Shifts Driving Housing Demand
% The Math of Household Unbundling (Atomización Familiar):

% In the 1960s–70s (~3.4–3.8 people/household): To house Santa Cruz’s ~210,000 residents, the city needed around 61,000 dwellings.
% Today (~2.2 people/household, according to ISTAC): Housing those same ~210,000 residents requires more than 95,000 dwellings.
% Result: Even without population growth, demographic fragmentation alone generates an objective deficit of over 34,000 additional housing units.
% Hyper-Motorization and the "Garage Deficit":

% Tenerife has one of the highest motorization rates in Europe, with over 830–840 vehicles per 1,000 inhabitants (ISTAC). Nearly every adult in a working household owns or requires a vehicle.
% The Legacy Problem: Santa Cruz’s 20th-century housing stock (neighborhoods like Ofra, Los Gladiolos, La Salud, or García Escámez built between 1960 and 1980) was constructed with zero underground garages, assuming street-level parking for a largely carless working-class population.
% Today’s Requirement: Modern households consider at least one (often two) private parking spaces non-negotiable. Because older consolidated buildings cannot retrofit underground basements into steep volcanic topography, new demand concentrates exclusively on modern developments.
% Evolving Standards of Living Space (Superficie y Confort):

% Mid-century social housing (viviendas de absorción) crammed three bedrooms and a bathroom into 50–65 m².
% Contemporary emancipated households demand larger living spaces, dedicated work/study rooms for teleworking, storage rooms (trasteros), accessibility (elevators), and private terraces or balconies.
% The Territorial Consequence: Because the paralysis of the PGO prevents Santa Cruz from developing new residential sectors with modern standards and garages, younger generations are priced out or forced to migrate to peripheral municipalities (La Laguna, El Rosario, Candelaria), creating the daily traffic congestion on the TF-5 and TF-1 highways (las colas de la autopista).

% La necesidad de vivienda se ve agravada por el encogimiento del tamaño de los hogares: en 1960 cada vivienda ocupaba, en promedio, 3,4 personas, mientras que en la actualidad el promedio se sitúa en 2,2 personas, lo que exige disponer de más unidades para acoger a una población constante

La presión sobre el mercado inmobiliario de Santa Cruz de Tenerife se ve fundamentalmente amplificada por tres transformaciones estructurales de carácter cualitativo e intraurbano:

\begin{enumerate}
  \item \textbf{Atomización familiar y encogimiento del tamaño del hogar}: Si hacia 1960 cada vivienda albergaba un promedio de 3,4 personas, en la actualidad el tamaño medio del hogar en el municipio se sitúa en torno a las 2,2 personas (ISTAC). Esta fragmentación significa que, para alojar a un volumen de población prácticamente estacionario (en torno a los 210.000 habitantes), la ciudad requiere hoy más de 34.000 unidades residenciales adicionales respecto a las que bastaban hace cuatro décadas, impulsada por el aumento de hogares unipersonales y la emancipación de adultos jóvenes.

  \item \textbf{Hipermotorización y necesidad crítica de aparcamiento}: Con una ratio insular que supera los 830 vehículos por cada 1.000 habitantes, Tenerife ostenta una de las tasas de motorización más elevadas de Europa. Gran parte del parque edificado consolidado en la segunda mitad del siglo XX (barriadas de los años 60 y 70 como Ofra, Los Gladiolos o La Salud) fue levantado sin plazas de garaje subterráneas. Para las nuevas unidades familiares, la vinculación obligatoria de una o dos plazas de aparcamiento por vivienda constituye una exigencia funcional ineludible que la trama consolidada no puede resolver, trasladando la presión hacia nuevos desarrollos con sótanos dotacionales.

  \item \textbf{Evolución de los estándares de superficie y habitabilidad}: Los requerimientos espaciales de la sociedad actual han dejado obsoleta la tipología de vivienda mínima de absorción (inmuebles de 50 a 65 $m^2$ para familias numerosas). La demanda contemporánea exige estancias de mayores dimensiones, espacios polivalentes adaptados al teletrabajo, cuartos trasteros y accesibilidad universal.
\end{enumerate}

\begin{figure}[htbp]
  \centering
  \includegraphics[width=0.95\linewidth]{figures/income_trend.pdf}
  \caption{Evolución de la Renta Media en Santa Cruz de Tenerife (2011-2021). Fuente: Elaboración propia a partir de datos del ISTAC e INE.}
  \label{fig:renta}
\end{figure}

Al encontrarse el planeamiento municipal bloqueado y sus sectores de expansión paralizados, la ciudad no ha podido ofertar esta tipología de vivienda moderna y dotada de garajes, expulsando a parte de la demanda emancipada hacia la periferia metropolitana y retroalimentando el colapso viario diario de acceso a la capital a través de las autopistas TF-5 y TF-1.
